\changecolours{ScoutGreen}
\section[Laboratório em Python]{Laboratório Computacional: Validação de Relacionamentos em Python}

% ------------------------------------------------------------------------------
\twocolframe[title={O Metamodelo Conceitual em Python},leftcol=7.1cm,rightcol=6.9cm]{%
  \alert{Por Que Simular o Modelo em Python?}\par
  \itemseps[topsepi=1mm,itemsepi=1.5mm,labelsep=2mm,leftmargini=3.5mm]
  \begin{itemize}
    \item \textbf{Validação Semântica Antes da Tecnologia:} Permite testar as regras de negócio em memória, garantindo que o modelo conceitual seja consistente.
    \item \textbf{Sem Dependência de SQL:} Nesta fase, nosso foco é 100\% nas regras formais do MER (Peter Chen, 1976).
    \item \textbf{Cenários Testados:}\\
      1. Regra 1:1 (\texttt{GERENCIA}): Impede que um departamento tenha dois gerentes.\\
      2. Regra 1:N (\texttt{LOTADO\_EM}): Impede que um funcionário pertença a dois departamentos.\\
      3. Regra N:M (\texttt{TRABALHA\_EM}): Registra horas semanais.
  \end{itemize}
}{%
  \alert{Objetivos do Laboratório em Memória}\par
  \itemseps[topsepi=1mm,itemsepi=1.5mm,labelsep=2mm,leftmargini=3.5mm]
  \begin{itemize}
    \item \textbf{Abstração Pura:}
      Compreender restrições conceituais sem o ruído sintático de dialetos de SGBDs específicos.
    \item \textbf{Tratamento de Exceções:}
      Disparar erros explícitos (\texttt{ValueError}) no momento em que operações violarem regras 1:1 ou 1:N.
    \item \textbf{Auditoria de Integridade:}
      Função de varredura que detecta instâncias órfãs que violam a participação total obrigatória.
    \item \textbf{Fixação Algorítmica:}
      Aproximar o raciocínio matemático de conjuntos da lógica de estruturas de dados em Python.
  \end{itemize}
}

% ------------------------------------------------------------------------------
\begin{frame}[fragile]{Código Python: Estrutura da Classe do Relacionamento}
  \vspace{0.1cm}
  \begin{columns}[t]
    \begin{column}{0.48\textwidth}
      \footnotesize
      \alert{Representação em Memória}\par
      \vspace{0.15cm}
      \begin{itemize}\setlength{\itemsep}{1.5mm}\setlength{\parskip}{0pt}
        \item Definimos enumerações para cardinalidade (\texttt{1:1}, \texttt{1:N}, \texttt{N:M}) e participação (\texttt{TOTAL}, \texttt{PARCIAL}).
        \item A classe gerencia o conjunto de instâncias e aplica as travas de cardinalidade.
      \end{itemize}
    \end{column}

    \begin{column}{0.48\textwidth}
      \alert{Definição em Python}\par
      \vspace{0.1cm}
\begin{lstlisting}[language=Python,basicstyle=\ttfamily\fontsize{5.2}{6.2}\selectfont]
class DefinicaoRelacionamento:
    def __init__(self, nome, ent_a, ent_b, 
                 cardinalidade, part_a, part_b,
                 papeis=None, atributos=None):
        self.nome = nome
        self.ent_a = ent_a
        self.ent_b = ent_b
        self.cardinalidade = cardinalidade
        self.part_a = part_a
        self.part_b = part_b
        self.papeis = papeis
        self.atributos = atributos or []
        self.instancias = []
\end{lstlisting}
    \end{column}
  \end{columns}
\end{frame}

% ------------------------------------------------------------------------------
\begin{frame}[fragile]{Código Python: Bloqueio Semântico de Violações}
  \vspace{0.1cm}
  \begin{columns}[t]
    \begin{column}{0.48\textwidth}
      \footnotesize
      \alert{Validação de Cardinalidade 1:1 e 1:N}\par
      \vspace{0.15cm}
      \begin{itemize}\setlength{\itemsep}{1.5mm}\setlength{\parskip}{0pt}
        \item \textbf{Em 1:1:} Se qualquer uma das entidades já estiver em uma instância, o método rejeita uma nova ligação.
        \item \textbf{Em 1:N:} Se a entidade do lado filho já estiver ligada a um pai, ela não pode ser cadastrada em outro.
      \end{itemize}
    \end{column}

    \begin{column}{0.48\textwidth}
      \alert{Método adicionar\_ligacao}\par
      \vspace{0.1cm}
\begin{lstlisting}[language=Python,basicstyle=\ttfamily\fontsize{5.2}{6.2}\selectfont]
def adicionar_ligacao(self, id_a, id_b, attrs=None):
    # Validacao de Cardinalidade 1:1
    if self.cardinalidade == TipoCardinalidade.UM_PARA_UM:
        for inst in self.instancias:
            if inst.id_a == id_a or inst.id_b == id_b:
                raise ValueError("Violacao de exclusividade 1:1!")

    # Validacao de Cardinalidade 1:N
    elif self.cardinalidade == TipoCardinalidade.UM_PARA_MUITOS:
        for inst in self.instancias:
            if inst.id_b == id_b:
                raise ValueError(f"[{id_b}] ja pertence a [{inst.id_a}]!")

    self.instancias.append(InstanciaRelacionamento(id_a, id_b, attrs))
\end{lstlisting}
    \end{column}
  \end{columns}
\end{frame}
